\DocumentMetadata{
  pdfversion=2.0,
  pdfstandard=ua-2,
  lang=en-US,
  tagging=on,
  tagging-setup={math/setup=mathml-SE,tabsorder=structure}
}
\documentclass[11pt,letterpaper]{article}
\usepackage[margin=0.68in,includefoot]{geometry}
\usepackage{newcomputermodern}
\usepackage{amsmath,amscd,mathtools}
\usepackage{tikz,adjustbox}
\providecommand{\square}{\mdlgwhtsquare}
\usepackage[hidelinks]{hyperref}
\hypersetup{
  pdftitle={Algebra PhD Exam, August 2016},
  pdfauthor={University of Florida Department of Mathematics},
  pdfsubject={Graduate mathematics examination},
  pdfdisplaydoctitle=true,
  bookmarksopen=true
}
\setcounter{secnumdepth}{0}
\setlength{\parindent}{0pt}
\setlength{\parskip}{0.28em}
\setlength{\emergencystretch}{3em}
\setlength{\leftmargini}{2.15em}
\setlength{\leftmarginii}{2.65em}
\setlength{\leftmarginiii}{3em}
\pagestyle{empty}
\raggedbottom
\allowdisplaybreaks
\NewTaggingSocketPlug{block/list/label}{examproblem}{%
  \tagstructbegin{tag=Lbl}\tagstructend%
  \tagstructbegin{tag=\LBody}%
  \tagstructbegin{tag=\ProblemHeadingTag,title-o={\ProblemHeadingText}}%
  \tagmcbegin{tag=\ProblemHeadingTag,actualtext={\ProblemHeadingText}}%
  #2%
  \tagmcend%
  \tagstructend%
}
\newenvironment{examproblems}[2]{%
  \def\ProblemHeadingTag{#2}%
  \AssignTaggingSocketPlug{block/list/label}{examproblem}%
  \begin{enumerate}%
  \setcounter{enumi}{#1}%
  \renewcommand{\labelenumi}{\textbf{\theenumi.}}%
  \setlength{\topsep}{0.35em}%
  \setlength{\partopsep}{0pt}%
  \setlength{\itemsep}{0.48em}%
  \setlength{\parsep}{0.18em}%
}{\end{enumerate}}
\newenvironment{examsubparts}{%
  \AssignTaggingSocketPlug{block/list/label}{default}%
  \begin{enumerate}%
  \setlength{\topsep}{0.18em}%
  \setlength{\partopsep}{0pt}%
  \setlength{\itemsep}{0.16em}%
  \setlength{\parsep}{0.1em}%
}{\end{enumerate}}
\newenvironment{examinstructions}{%
  \par\begingroup\itshape
}{%
  \par\endgroup\vspace{0.18em}
}
\makeatletter
\renewcommand\subsection{\@startsection{subsection}{2}{0pt}%
  {-1.25ex plus -.25ex minus -.1ex}%
  {0.55ex plus .1ex}%
  {\normalfont\large\bfseries}}
\renewcommand{\@maketitle}{%
  \newpage\null\vskip -1em%
  {\footnotesize\bfseries
    \href{https://www.ufl.edu/}{University of Florida}, \href{https://math.ufl.edu/}{Department of Mathematics}\par}%
  \vskip -0.5em%
  \tagstructbegin{tag=H1,title={Algebra PhD Exam, August 2016}}%
  \tagpdfparaOff%
  \tagmcbegin{tag=H1}%
    {\LARGE\bfseries\href{https://gma.math.ufl.edu/exams/algebra-phd/}{Algebra PhD Exam}, August 2016\par}%
  \tagmcend%
  \tagpdfparaOn%
  \tagstructend%
  \vskip 0.25em%
  \tagmcbegin{artifact}%
    \hrule height 1pt%
  \tagmcend%
  \vskip 1em%
}
\makeatother
\title{Algebra PhD Exam, August 2016}
\author{}
\date{}
\begin{document}
\maketitle
\thispagestyle{empty}
\begin{examinstructions}
Answer seven problems. Write your answers clearly in complete English sentences. You may quote results (within reason) as long as you state them clearly.
\end{examinstructions}
\begin{examproblems}{0}{H2}
\def\ProblemHeadingText{Problem 1.}
\item
Let~\(L\) be a splitting field over~\(\mathbb{Q}\) for the polynomial \(f(X)=X^3-7\).
\begin{examsubparts}
\item[\textnormal{(a)}]
Determine the Galois group \(\operatorname{Gal}(L/\mathbb{Q})\).
\item[\textnormal{(b)}]
For each subgroup~\(H\) of \(\operatorname{Gal}(L/\mathbb{Q})\) determine the fixed field \(L^H\) of~\(H\).
\end{examsubparts}
\def\ProblemHeadingText{Problem 2.}
\item
Let~\(p\) be prime and let \(n\geq 1\).
\begin{examsubparts}
\item[\textnormal{(a)}]
Prove that there is a field \(\mathbb{F}_{p^n}\) with \(p^n\) elements.
\item[\textnormal{(b)}]
Prove that if~\(K\) is a field with \(p^n\) elements then \(K\cong\mathbb{F}_{p^n}\).
\end{examsubparts}
\def\ProblemHeadingText{Problem 3.}
\item
This problem has three parts.
\begin{examsubparts}
\item[\textnormal{(a)}]
Let \(\mathcal C\) be a category. Define what it means for a \(\mathcal C\)-morphism \(f:X\to Y\) to be a monomorphism.
\item[\textnormal{(b)}]
Give the definition of a concrete category.
\item[\textnormal{(c)}]
Let \(\mathcal C\) be a concrete category and let \(f:X\to Y\) be a \(\mathcal C\)-morphism whose underlying set map is one-to-one. Prove that~\(f\) is a monomorphism.
\end{examsubparts}
\def\ProblemHeadingText{Problem 4.}
\item
Prove that the group~\(G\) defined by the generators and relations
\[
G=\langle x,y\mid xyx=y\rangle
\]
 is infinite and nonabelian.
\def\ProblemHeadingText{Problem 5.}
\item
Let~\(R\) be a commutative ring with~\(1\) and let \(I,J\) be ideals in~\(R\). Prove that there is an isomorphism of~\(R\)-modules
\[
(R/I)\otimes_R(R/J)\cong R/(I+J).
\]
\def\ProblemHeadingText{Problem 6.}
\item
Let~\(R\) be a commutative ring with~\(1\).
\begin{examsubparts}
\item[\textnormal{(a)}]
Define what it means for a left~\(R\)-module to be flat.
\item[\textnormal{(b)}]
Prove that if~\(F\) is a free left~\(R\)-module then~\(F\) is flat.
\end{examsubparts}
\def\ProblemHeadingText{Problem 7.}
\item
Let \(R\subset S\) be integral domains such that~\(S\) is integral over~\(R\). Prove that~\(R\) is a field if and only if~\(S\) is a field.
\def\ProblemHeadingText{Problem 8.}
\item
Let~\(R\) be a commutative ring with~\(1\).
\begin{examsubparts}
\item[\textnormal{(a)}]
Define the Jacobson radical \(J(R)\) of~\(R\).
\item[\textnormal{(b)}]
Let~\(M\) be a finitely generated~\(R\)-module and let~\(I\) be an ideal of~\(R\) which is contained in \(J(R)\). Prove that if \(IM=M\) then \(M=\{0\}\).
\end{examsubparts}
\def\ProblemHeadingText{Problem 9.}
\item
Determine which of the following rings are Dedekind domains, with some explanation in each case.
\begin{examsubparts}
\item[\textnormal{(a)}]
\(\mathbb{F}_7[X]\)
\item[\textnormal{(b)}]
\(\mathbb{F}_7[X,Y]\)
\item[\textnormal{(c)}]
\(\mathbb{Z}[\sqrt{3}]\)
\item[\textnormal{(d)}]
\(\mathbb{Z}[\sqrt{5}]\)
\end{examsubparts}
\def\ProblemHeadingText{Problem 10.}
\item
Let~\(G\) be a finite group and let~\(K\) be a field of characteristic~\(0\). Let~\(V\) be a \(K[G]\)-module which is a finite-dimensional vector space over~\(K\) and let~\(W\) be a \(K[G]\)-submodule of~\(V\). Prove there is a \(K[G]\)-submodule \(W'\) of~\(V\) such that \(V=W\oplus W'\).
\def\ProblemHeadingText{Problem 11.}
\item
Determine the number of isomorphism classes of semisimple rings with \(4000\) elements. You don’t have to list the rings, but you do need to justify your computations.
\end{examproblems}
\end{document}
